\BeleskeID{09_2-s016}
% element: 09_2-s016:e04
Kolo je neopterećeno, pa zakočeni N tranzistor ne odvodi struju sa izlaza. P tranzistor mora imati nultu struju iako ima uslove za formiranje kanala. To se ostvaruje pri nultom naponu između drejna i sorsa u omskoj oblasti.

Pretpostavka zasićenja dala bi
\[I_{Dp}=\frac{k_p}2(V_{GS,P}-V_{Tp})^2=\frac{k_p}2(-V_{DD}-V_{Tp})^2>0\]
Ta pozitivna struja nije saglasna sa ravnotežom struja neopterećenog izlaza. Zato izlaz dostiže $V_{DD}$. Negativan prag P tranzistora treba imati u vidu pri proveri uslova provođenja.
